\chapter{AI工具箱：理解你的协作伙伴}

第一次使用生成式 AI 时，人们往往会被它流畅的表达和快速的执行能力吸引。它可以在几秒钟内解释一个概念，生成数十行代码，整理一份数据分析计划，甚至给出一篇结构完整的报告。这样的体验很容易让人产生两种相反但同样危险的判断：一种是认为 AI 无所不能，可以直接替代专业人员；另一种是发现 AI 会犯错后，便认为它不值得使用。

更准确的认识介于两者之间。AI 既不是能够独立承担责任的专家，也不只是一个更快的搜索框。它是一种能够理解自然语言指令、生成候选方案并调用其他工具完成任务的通用协作系统。它的价值取决于三个条件：任务是否适合交给 AI，使用者是否提供了足够的信息，以及结果是否经过了与风险相匹配的验证。

本章的目标，是建立一套关于 AI 协作的基本心智模型。学完本章后，读者应当能够回答以下问题：

\begin{itemize}
  \item 大语言模型为什么能够完成看似不同的任务？
  \item AI 擅长什么，又为什么会稳定地产生某些错误？
  \item 面对一个真实任务，应该选择哪一类 AI 工具？
  \item 哪些信息可以交给 AI，哪些判断必须由人负责？
  \item 如何把一次随意问答升级为可检查、可复现的协作过程？
\end{itemize}

% ----------------------------------------------------------------
\section{AI协作系统是如何工作的}

今天人们所说的“AI 工具”，通常不是一个孤立模型，而是由基础模型、系统提示、检索模块、代码执行环境、外部工具和权限控制共同组成的协作系统。基础模型决定能力上限，系统设计决定能力以何种方式被调用，使用者提供的上下文则决定当前任务究竟被解释成什么。将这三个层次混在一起，会使人把一次偶然成功误认为稳定能力，也会把界面限制误认为模型本身的限制。基础模型研究因此强调，通用能力能够迁移到大量下游任务，但这种广泛适用性同时带来难以预先穷举的失效方式与社会风险\cite{bommasani2021opportunities}。
% ----------------------------------------------------------------

\subsection{大语言模型的基本工作方式}

现代大语言模型通常建立在 Transformer 架构之上，通过注意力机制在上下文中建立词元之间的关系\cite{vaswani2017attention}。训练目标使模型善于根据已有上下文预测后续内容；规模扩大以后，这种预测能力表现为解释、摘要、翻译、代码生成和少样本任务适应等多种外在能力\cite{brown2020language}。但这里必须保留一个重要区分：语言上连贯的回答，是训练目标直接鼓励的结果；事实正确、因果成立和方案可执行，则需要额外证据与验证。

当前许多生成式 AI 工具的核心是大语言模型。它接收一段上下文，并根据上下文生成接下来最合适的内容。这里的“内容”可以是自然语言，也可以是代码、公式、表格或结构化数据。模型在大量文本和代码上学习了语言模式、概念联系与常见问题的解决形式，因此能够在新的任务中组合已有知识，生成看起来具有推理能力的回答。

但是，大语言模型生成回答的过程与人类查阅权威数据库不同。它并不是先找到一条已经被验证的事实，再原样返回；它更像是在当前上下文中，逐步生成最可能且最符合指令的表达。因此，语言流畅、结构完整和语气确定，并不能证明内容正确。

理解这一点非常重要。大语言模型的输出首先是一个\textbf{候选答案}，而不是自动成立的事实。候选答案可能非常有价值，也可能混合正确内容、合理猜测和完全错误的信息。使用者的任务，是决定它需要接受何种检查之后，才能被采用。

\subsection{从基础模型到AI协作系统}

实际使用的 AI 工具通常不只是一个孤立的大语言模型。现代 AI 系统可能在模型之外增加多种能力：

\begin{description}
  \item[上下文管理] 保存当前对话、项目文件和用户提供的背景信息，使模型能够围绕同一个任务持续工作。
  \item[信息检索] 从网页、文献库、组织内部知识库或项目文档中查找资料，再结合资料生成回答。
  \item[代码执行] 在受控环境中运行代码、读取数据、生成图表并返回运行结果。
  \item[工具调用] 调用搜索、数据库、终端、地图服务、软件接口或其他外部工具。
  \item[文件编辑] 读取多个项目文件，修改代码，并通过测试或编译检查结果。
\end{description}

因此，当我们评价一个 AI 工具时，不能只问“它使用了什么模型”，还要问：它能够访问什么上下文，能够调用什么工具，执行过程是否可见，生成结果是否可以验证，以及数据会被发送到哪里。

\subsection{为什么它既聪明又不可靠}

这种矛盾并不是偶然缺陷，而是生成式模型工作方式的直接结果。模型可以组合训练中见过的模式，却不天然拥有对现实世界的持续观测，也不会自动标记每句话的证据来源。相关研究用“随机鹦鹉”的比喻提醒使用者：流畅文本容易诱发理解和权威感，但语言形式本身不能证明模型真正掌握了所讨论对象\cite{bender2021stochastic}。因此，评价 AI 输出时必须把\textbf{表达质量}、\textbf{任务适切性}与\textbf{事实可靠性}分开判断。

AI 的能力与局限并不矛盾，它们来自同一个根源。理解这一点，有助于建立合理预期，避免因为 AI 的流畅表达而高估其可靠性。

\paragraph{擅长模式组合，不保证事实成立}

大语言模型从大量材料中学习了问题与回答之间的模式。遇到常见任务时，它能够迅速组合出高质量方案；遇到信息不足、专业长尾或新颖任务时，它也会继续生成形式合理的回答，而不是自动停止。可以把它想象为一个读过海量资料、能够熟练套用模式的助手：遇到常见题型，它给出高质量答案；遇到需要真实核查的细节，它仍然会按“最符合语境”的方式继续作答，而不会自动承认不确定。

这意味着，任务越接近广泛存在且容易验证的模式，AI 通常越可靠；任务越依赖稀有知识、内部信息或复杂责任判断，越需要谨慎。

\paragraph{回答高度依赖上下文}

AI 的输出取决于当前获得的上下文。同一个问题，如果补充了数据字典、目标定义和错误成本，回答质量可能显著提高。反过来，如果对话过长、信息相互冲突或关键材料没有被提供，AI 可能忽略重要条件。

因此，使用者需要主动管理上下文：

\begin{itemize}
  \item 只提供与当前任务相关且必要的信息；
  \item 明确指出哪些材料具有最高优先级；
  \item 对关键定义和约束进行重复确认；
  \item 将复杂项目拆成可检查的阶段；
  \item 在阶段结束时记录已确认的决策。
\end{itemize}

\paragraph{输出具有不稳定性}

对于开放性任务，AI 在相同输入下可能生成不同答案。这种变化有时有益，因为它能够提供更多思路；但对于需要可复现的研究与工程流程，不稳定性会带来问题。

要提高可复现性，应保存使用的输入材料、关键 Prompt、模型或工具版本、生成代码、依赖环境和验证结果。对于核心步骤，不应只保存最后一段文字结论，而应保存能够重新运行和检查的过程。

% ----------------------------------------------------------------
\section{AI的能力边界：能做什么，不能做什么}
% ----------------------------------------------------------------

理解 AI 的价值，需要同时看清它擅长的任务和它稳定失效的场景。两者共同构成协作决策的判断基础：前者告诉你在哪里使用 AI，后者告诉你在哪里必须亲自把关。

\subsection{AI擅长的五类任务}

AI 最擅长的任务通常具有一个共同特点：使用者能够描述目标，结果可以通过阅读、运行、比较或测试进行检查。AI 可以显著降低这些任务的启动成本，并快速生成多个可供选择的方案。

\paragraph{解释与学习：把陌生知识转化为可理解的入口}

AI 可以用不同深度解释同一个概念，比较相近方法，补充学习路径，并根据使用者的背景调整表达。例如，在学习交通流预测时，可以先让 AI 用直观语言解释时间序列中的趋势、周期和滞后，再要求它给出数学定义和代码示例。

这种能力适合建立初步理解，但不适合替代教材、论文和标准。一个可靠的学习过程通常是：

\begin{enumerate}
  \item 让 AI 帮助识别概念、关键词和知识结构；
  \item 使用教材、原始论文或官方文档核实关键内容；
  \item 回到 AI，针对仍不理解的部分继续提问；
  \item 通过练习、推导或实验检验自己是否真正理解。
\end{enumerate}

AI 在这里扮演的是导师助理，而不是唯一知识来源。

\paragraph{生成与修改代码：降低实现成本}

代码生成是 AI 最直接的价值之一。它可以根据需求生成初始代码，解释已有程序，补充注释，定位报错，编写测试，或把重复操作封装为函数。对于熟悉问题但不熟悉某个软件库的使用者，AI 尤其能够减少查找接口和样板代码的时间。

然而，判断一段 AI 生成代码是否可用，至少需要检查四个层面：

\begin{enumerate}
  \item \textbf{能否运行：}语法是否正确，依赖是否存在，输入输出是否匹配；
  \item \textbf{是否按要求运行：}代码执行的任务是否与需求一致；
  \item \textbf{方法是否合理：}数据处理、模型选择和评价方式是否正确；
  \item \textbf{结果是否可信：}输出能否通过测试、基准和领域知识检验。
\end{enumerate}

“代码成功运行”只回答了第一个问题。真实项目中的严重错误，往往发生在后三个层面。

\paragraph{数据分析：加速探索，但不替代数据理解}

在数据分析任务中，AI 可以帮助生成描述统计、缺失值检查、分布图、相关性分析和初步建模代码，也可以根据输出结果提出下一步探索建议。它适合承担大量机械性工作，让使用者把精力集中在异常现象、业务含义和分析决策上。

例如，面对道路检测器数据，AI 可以快速计算每个检测器的缺失率，并绘制速度随时间变化的曲线。但它不知道某段缺失是设备故障、通信中断还是道路封闭，也不知道速度为零是否代表拥堵、停车或错误记录。这些解释依赖数据采集机制和领域知识。

因此，AI 可以帮助回答“数据呈现了什么模式”，但“这些模式意味着什么”仍需要人与证据共同回答。

\paragraph{方案生成与比较：扩展人的搜索空间}

AI 可以在短时间内生成多个候选方案，并按成本、风险、适用条件和验证难度进行比较。例如，对于“预测公交站点下一小时客流”的任务，可以要求 AI 分别提出简单基准、传统机器学习方法和时序神经网络方法，再比较每种方法所需的数据、解释性和部署成本。

在这一类任务中，AI 最有价值的地方不是替人决定，而是帮助人避免过早锁定第一个想法。一个成熟的使用方式是要求 AI：

\begin{itemize}
  \item 给出至少两种性质不同的方案；
  \item 说明每种方案依赖的假设；
  \item 列出方案可能失败的原因；
  \item 提出区分这些方案的实验；
  \item 明确仍需要人决定的问题。
\end{itemize}

\paragraph{表达与沟通：帮助组织已有认识}

AI 可以协助整理报告结构、改写说明文字、生成图表标题、提炼会议记录，并把技术分析转换为不同受众能够理解的表达。对于需要同时面向研究人员、管理者和公众的工程项目，这种转换能力十分实用。

但表达能力也可能掩盖内容缺陷。AI 可以把证据不足的结论写得非常完整。因此，在使用 AI 改写报告时，应先确认事实、数字和结论，再优化语言。不要让更流畅的表达提高一个错误结论的说服力。

\subsection{AI稳定失效的场景}

理解 AI 的局限，不是为了减少使用，而是为了决定应当在哪里增加检查、降低权限或改由人处理。

\paragraph{它不会自动理解你的真实问题}

现实问题往往包含没有被写出的目标和约束。使用者说“建立一个最准确的事故预测模型”，AI 可能自然地追求某个测试集指标；但真实目标也许是尽量减少严重事故漏报，或帮助有限的巡查资源确定重点区域。“最准确”并不是一个完整的问题定义。

AI 只能根据被提供的信息工作。当目标含糊时，它通常会用常见假设填补空白。这些假设可能合理，也可能完全偏离真实需求。因此，问题定义不能外包给 AI。AI 可以帮助澄清问题，但最终目标必须由任务负责人确认。

\paragraph{它不会自动理解数据如何产生}

数据不是现实的完整复制，而是某种采集过程留下的记录。同样一个数值，在不同采集机制下可能有完全不同的含义。

以速度为零为例：

\begin{itemize}
  \item 在 GPS 轨迹中，它可能代表车辆真实停止；
  \item 在道路检测器数据中，它可能代表设备故障或未检测到车辆；
  \item 在经过聚合的数据中，它可能是缺失值被错误填充后的结果；
  \item 在仿真输出中，它可能代表模型中的特殊状态。
\end{itemize}

如果没有数据字典、采集说明和场景知识，AI 无法仅凭列名可靠地区分这些情况。让 AI 自动清洗数据之前，必须先解释数据来源和字段语义。

\paragraph{它不知道未被提供的隐含约束}

专业工作中存在大量隐含约束。例如，交通预测不能使用预测时刻之后才能获得的信息；政策评估不能把相关性直接解释为因果关系；包含个人轨迹的数据不能随意上传到公共服务；安全相关决策不能只依据一个无法解释的模型输出。

领域专家常常因为长期经验而默认这些约束显而易见，但对 AI 来说，未被表达的约束通常等于不存在。可靠协作要求使用者把关键隐含知识转化为明确规则，并在结果中检查这些规则是否被遵守。

\paragraph{它可能生成不存在或不准确的信息}

大语言模型可能编造文献、软件参数、法规条款、数据来源或计算结果，这类现象通常被称为幻觉。幻觉并不总是表现为荒谬内容。更危险的情况是，回答大部分正确，只在一个关键数字、引用或条件上出错。

对于可核实事实，应优先使用原始来源：

\begin{itemize}
  \item 文献结论应查阅原始论文；
  \item 软件接口应查阅官方文档并实际运行；
  \item 法规与标准应查阅正式发布版本；
  \item 数据集字段应查阅数据字典；
  \item 计算结果应由代码或独立方法复算。
\end{itemize}

\paragraph{它不能承担专业与伦理责任}

AI 不会因错误分析承担工程事故、研究不端、隐私泄露或政策误判的责任。无论报告、代码或建议由谁生成，采用它们的人和组织仍然需要对结果负责。

因此，越是高风险的任务，AI 越不应该直接给出最终决定。它可以整理证据、生成候选方案和提示风险，但安全判断、政策取舍、利益协调和责任承诺必须由有授权的人完成。

\subsection{三个不能混淆的判断}

以上能力与局限，可以归纳为三个在实际工作中容易被混淆的判断。它们贯穿本书始终，也是建立可靠协作的前提。

\begin{enumerate}
  \item \textbf{模型能力不等于事实可靠性。} 模型能够解释复杂概念，并不意味着它给出的每个事实、数字和引用都正确。
  \item \textbf{生成代码不等于解决问题。} 一段代码可以成功运行，却仍然使用了错误的数据、错误的方法或错误的评价标准。
  \item \textbf{自动化不等于自主负责。} AI 可以自动执行多个步骤，但任务目标、风险边界和最终责任仍然属于人。
\end{enumerate}

% ----------------------------------------------------------------
\section{AI工具生态与选择方法}
% ----------------------------------------------------------------

AI 工具更新很快，具体产品名称和功能会变化。比记住产品列表更重要的是理解不同工具的交互方式、上下文范围和风险边界。

\subsection{主要工具类别}

可以将常见 AI 工具分为以下五类：

\paragraph{对话型AI：适合探索、解释与方案讨论}

对话型 AI 以自然语言交流为主，适合概念解释、头脑风暴、任务拆解、文字修改和方案比较。这类工具最具代表性的有 OpenAI 的 ChatGPT、Anthropic 的 Claude 以及 Google 的 Gemini，它们通常以网页或应用程序界面的形式提供服务，无需本地配置即可开始使用（见图~\ref{fig:chat_ai_tools}）。它的优势是启动成本低，适合在问题尚未完全清楚时展开讨论。

\begin{figure}[H]
  \centering
  % 请将以下三行替换为实际截图或 logo 文件，例如：
  % \includegraphics[width=0.3\textwidth]{figures/chatgpt.png}
  % \includegraphics[width=0.3\textwidth]{figures/claude.png}
  % \includegraphics[width=0.3\textwidth]{figures/gemini.png}
  \fbox{\parbox{0.28\textwidth}{\centering\vspace{1.5cm}ChatGPT\vspace{1.5cm}}}
  \hfill
  \fbox{\parbox{0.28\textwidth}{\centering\vspace{1.5cm}Claude\vspace{1.5cm}}}
  \hfill
  \fbox{\parbox{0.28\textwidth}{\centering\vspace{1.5cm}Gemini\vspace{1.5cm}}}
  \caption{主流对话型 AI 工具（从左至右：ChatGPT、Claude、Gemini）}
  \label{fig:chat_ai_tools}
\end{figure}

它的局限是对真实项目上下文了解有限。如果使用者只粘贴少量代码或局部数据，AI 可能在缺少整体结构的情况下提出修改建议。对话型 AI 更适合回答”我应该考虑什么”和”有哪些可能方案”，不宜在缺少项目材料时直接承担大范围修改。

\paragraph{代码协作型AI：适合在真实项目中执行}

代码协作型 AI 可以读取项目文件、搜索代码、修改内容、运行命令并检查测试。它适合跨文件开发、错误修复、重构、批量修改和构建自动化流程。代表性工具包括 Anthropic 的 Claude Code（可在终端中直接与整个代码库交互）、GitHub Copilot（集成在 VS Code 等编辑器中，提供代码补全与内联对话）以及 Cursor（内置 AI 能力的代码编辑器）。OpenAI 早期推出的 Codex 奠定了代码生成能力的基础，后续许多工具都在此基础上发展而来。与只看到一段粘贴代码的对话工具相比，代码协作型 AI 能够利用更完整的项目上下文（见图~\ref{fig:code_ai_tools}）。

\begin{figure}[H]
  \centering
  % 请将以下两行替换为实际截图文件，例如：
  % \includegraphics[width=0.45\textwidth]{figures/claude_code.png}
  % \includegraphics[width=0.45\textwidth]{figures/github_copilot.png}
  \fbox{\parbox{0.45\textwidth}{\centering\vspace{1.5cm}Claude Code\vspace{1.5cm}}}
  \hfill
  \fbox{\parbox{0.45\textwidth}{\centering\vspace{1.5cm}GitHub Copilot\vspace{1.5cm}}}
  \caption{代码协作型 AI 工具示例（Claude Code 与 GitHub Copilot）}
  \label{fig:code_ai_tools}
\end{figure}

代码协作型 AI 的权限也更高，因此需要更清楚的边界。使用前应确认：

\begin{itemize}
  \item 它可以读取和修改哪些目录；
  \item 它能否运行终端命令或访问网络；
  \item 修改前后是否能够查看差异；
  \item 是否存在测试、编译或其他验证手段；
  \item 关键修改是否需要人工确认。
\end{itemize}

\paragraph{数据分析型AI：适合快速探索和原型验证}

数据分析型 AI 通常能够接收文件、运行代码并生成图表，适合快速完成探索性数据分析、简单统计计算和原型验证。ChatGPT 的 Advanced Data Analysis 功能（早期称为 Code Interpreter）是这类工具的典型代表，用户可以直接上传数据文件，由 AI 在沙盒环境中运行 Python 代码并返回图表与统计结果。对于尚未配置本地环境的使用者，这类工具能够迅速展示数据中可能存在的模式。

但在上传数据之前，必须确认隐私、授权和保密要求。即使平台提供代码执行能力，也要检查它实际读取了哪些列、进行了哪些转换，以及图表和结论是否与原始数据一致。正式项目中，快速探索的结果通常还需要迁移到可复现的本地或组织内部流程。

\paragraph{检索与研究型AI：适合发现资料和整理证据}

检索与研究型 AI 可以帮助寻找关键词、筛选资料、比较观点和构建初步文献地图。Perplexity AI 是其中的代表，它将网络搜索与语言模型结合，在回答问题时附带可核查的来源链接。Elicit 和 Consensus 则专注于学术文献检索，能够快速梳理某一研究方向的主要发现与分歧点。这类工具适合在陌生领域中定位可能相关的材料。

然而，文献检索的目标不是生成一份看起来完整的引用列表，而是建立可核实的证据链。使用者应检查文献是否真实存在、是否真正支持相应结论、研究设计是否适用，以及资料是否已经过时。对于关键论断，应回到原始材料，而不是只依赖 AI 的摘要。

\paragraph{专业领域AI：适合受约束的特定工作流}

一些 AI 工具围绕特定领域设计，例如面向法律文书检索的 Harvey 和 CoCounsel，面向学术文献综述的 Consensus，以及部分医疗平台提供的 AI 辅助记录工具。这些工具往往集成了专业数据源、格式要求和工作流，因此可能比通用工具更适合特定任务。

“专业工具”并不自动意味着“结果正确”。选择时仍应检查数据来源、适用范围、更新频率、引用方式和责任边界。专业工具的价值在于缩小任务范围和提供领域上下文，而不是取消人工验证。

\subsection{根据任务特征选择工具}

选择工具时，不应只比较模型排行榜或宣传中的能力。真实项目更需要考虑任务性质、上下文位置、验证方式和数据风险。面对一个任务，可以依次提出以下问题：

\begin{enumerate}
  \item \textbf{需要什么结果？} 是解释、建议、代码、数据分析、文献证据，还是最终决策？
  \item \textbf{信息在哪里？} 在使用者脑中、一个文件中、整个代码库中、数据库中，还是公开资料中？
  \item \textbf{结果如何验证？} 能否运行测试、与基准比较、查阅原始来源或由专家复核？
  \item \textbf{错误代价多高？} 错误只是浪费几分钟，还是可能影响安全、隐私、研究结论或公共决策？
  \item \textbf{数据能否被提供？} 是否包含个人信息、商业秘密、未公开研究数据或受限制材料？
\end{enumerate}

这些问题决定了 AI 应扮演什么角色，也决定了工具需要具备什么能力。表~\ref{tab:tool_selection} 给出了一个根据任务特征选择工具的参考矩阵。在许多真实项目中，答案不是选择一个工具，而是组合使用多个：先使用对话型 AI 澄清任务，再使用代码协作型 AI 修改项目，最后使用独立测试和专业人员验证结果。

\begin{table}[H]
\centering
\caption{根据任务特征选择AI工具}
\label{tab:tool_selection}
\begin{tabular}{p{2.6cm}p{3.4cm}p{3.4cm}p{3.2cm}}
\toprule
任务类型 & 优先考虑的工具 & 主要优势 & 必须检查的风险 \\
\midrule
概念解释与方案讨论 & 对话型 AI & 启动快，适合迭代澄清 & 事实错误、遗漏约束 \\
项目代码修改 & 代码协作型 AI & 能读取项目并运行验证 & 修改范围、权限、回归错误 \\
快速数据探索 & 数据分析型 AI & 能运行代码并生成图表 & 隐私、字段误解、不可复现 \\
文献与资料发现 & 检索或研究型 AI & 能快速定位候选材料 & 虚假引用、证据不匹配 \\
高风险专业任务 & 专业工具加人工专家 & 领域资料和流程更完整 & 适用范围、责任与合规 \\
\bottomrule
\end{tabular}
\end{table}

\subsection{权限应与风险匹配}

一个工具能够读取数据、运行命令或生成决策建议，并不意味着它应该被允许这样做。权限应与任务风险匹配。

对于低风险、易恢复的任务，可以给予 AI 较高执行权限，例如整理公开资料或修改有完整测试覆盖的代码。对于涉及敏感数据、不可逆操作或高风险决策的任务，应限制访问范围，并要求人在关键节点确认。

% ----------------------------------------------------------------
\section{人机协作的基本框架}
% ----------------------------------------------------------------

将 AI 看作执行工程师，是一种实用但并不完全的类比。优秀的执行工程师可以快速实现方案、发现问题并提出建议，但项目负责人仍需定义目标、分配任务、提供资源、检查结果并承担责任。这一框架确立了人和 AI 在整个协作过程中的基本角色。

\subsection{人和AI的基本分工}

人机分工不是简单地把“重复劳动”交给 AI、把“创造性劳动”留给人。真实项目中的工作往往同时包含搜索、判断、实现和验证。更稳健的分法，是依据责任和可验证性分配任务：可快速检查、失败代价低的工作可以更多交给 AI；目标含糊、证据不足、失败代价高或涉及价值取舍的工作，应由人保持主导。这个分工原则也解释了为什么 AI 能显著降低生成代码和候选方案的成本，却不会自动降低系统集成、数据质量和长期维护的成本。机器学习系统中的隐性技术债务研究表明，模型代码常常只是完整系统的一小部分，数据依赖、配置、监控和反馈回路才是长期风险的重要来源\cite{sculley2015hidden}。

\begin{table}[H]
\centering
\caption{人和AI在协作项目中的基本分工}
\begin{tabular}{p{3.0cm}p{5.0cm}p{5.0cm}}
\toprule
工作环节 & 人的主要责任 & AI适合承担的工作 \\
\midrule
问题定义 & 确认真实目标、利益相关方与错误成本 & 提出澄清问题，帮助拆解任务 \\
方案设计 & 决定关键约束、评价标准与可接受风险 & 生成候选方案，比较优缺点 \\
技术实现 & 确认方法合理并控制修改范围 & 编写代码、整理数据、生成文档 \\
结果验证 & 设计验证流程并判断证据是否充分 & 运行检查、生成对比、提示潜在风险 \\
最终决策 & 承担专业、伦理与组织责任 & 整理证据和呈现备选方案 \\
\bottomrule
\end{tabular}
\end{table}

这一分工不是固定不变的。随着任务风险升高，AI 的角色应从执行者逐步降低为副驾驶或参谋；随着验证能力增强，一些重复、低风险步骤则可以逐步自动化。

\subsection{组织过程，而不只是索取答案}

较弱的 AI 使用方式是直接提问：“请告诉我应该使用哪个模型。”较强的方式是要求 AI 参与一个可检查的决策过程：

\begin{quote}
请先列出这个任务中仍不明确的信息，再提出三种复杂度不同的建模方案。对每种方案说明数据要求、关键假设、评价方法和可能失败的原因。暂时不要替我选择最终方案。
\end{quote}

后一种方式不会保证 AI 正确，但它暴露了更多中间信息，使使用者能够检查假设、比较方案并发现遗漏。与 AI 协作的质量，不只取决于最后得到什么答案，也取决于过程是否让关键决策变得可见。

\subsection{为AI设置停止条件}

项目负责人不仅要告诉 AI 做什么，也要告诉它何时应该停止并请求确认。典型停止条件包括：

\begin{itemize}
  \item 缺少完成任务所需的关键数据或定义；
  \item 多份材料之间存在无法自行解决的冲突；
  \item 操作可能删除、覆盖或公开重要内容；
  \item 结果涉及安全、政策、伦理或法律判断；
  \item 当前结果无法通过预先约定的验证标准；
  \item AI 对关键事实无法提供可核实来源。
\end{itemize}

允许 AI 表达不确定并停止，是可靠协作的重要条件。一个始终继续执行的系统，未必是一个值得信任的系统。

% ----------------------------------------------------------------
\section{可靠协作工作流与贯穿案例}
% ----------------------------------------------------------------

真正有价值的 AI 协作，不是偶然获得一次好答案，而是能够稳定完成一类任务。本节给出一套六步工作流，并通过贯穿案例展示它在真实项目中的应用方式。

\subsection{六步工作流}

\paragraph{第一步：定义任务与成功标准}

开始前先写清楚任务目标、交付物、范围和成功标准。例如，“分析交通数据”过于模糊，而“检查 20 个道路检测器在过去一个月中的缺失情况，并生成一份按检测器排序的缺失率表和每日缺失热力图”更容易执行和验证。

\paragraph{第二步：提供最小充分上下文}

上下文不是越多越好，而是应当足以支持当前任务。至少要提供与任务相关的数据定义、文件结构、已有约束和期望输出。对于不确定的信息，应明确标注，而不是让 AI 猜测。

\paragraph{第三步：让AI先提出计划}

对于多步骤任务，先让 AI 说明准备如何完成、需要读取哪些材料、会生成哪些结果以及准备如何验证。使用者确认计划后再执行，可以在成本较低时发现方向错误。

\paragraph{第四步：分阶段执行并保留中间结果}

将复杂任务拆成小步骤，每一步都有清楚输入和可检查输出。例如，先确认数据字段，再生成清洗规则；先建立简单基准，再尝试复杂模型。不要让 AI 在一个不可见的长过程里同时完成所有决策。

\paragraph{第五步：使用独立方法验证}

验证方式应尽可能独立于生成方式。AI 生成的代码可以通过测试和人工抽查验证；AI 汇总的事实可以通过原始来源验证；模型结果可以通过基准、留出数据和领域常识验证。不能只让同一次对话中的 AI 再次确认自己“是否正确”。

\paragraph{第六步：记录决策与责任}

保存关键 Prompt、输入材料、代码版本、运行环境、验证结果和人工决策。记录的目的不是保存所有对话，而是让其他人能够回答：做了什么，为什么这样做，哪些内容由 AI 生成，哪些内容经过了什么检查，以及谁批准了最终结果。

\subsection{贯穿案例：从任务提出到工具选择}

\begin{quote}
\textbf{【案例占位】}

这里将加入一个完整的交通运输工程案例，用于展示如何判断任务是否适合使用 AI、如何选择工具、如何划分人和 AI 的责任，以及如何建立可检查的协作流程。

\textbf{【待确定内容】} 真实任务背景、可用数据、候选 AI 工具、数据与权限风险、人机分工、验证方法、协作记录，以及采用 AI 前后工作方式的差异。
\end{quote}

案例最终需要体现本章的核心原则：工具选择必须服从任务目标；AI 可以负责生成候选答案和执行明确步骤，但问题定义、风险控制、结果验证与最终责任仍然属于人。

% ----------------------------------------------------------------
\section{开始项目前的AI协作检查清单}
% ----------------------------------------------------------------

在把任务交给 AI 之前，可以使用以下清单进行快速检查。

\begin{enumerate}
  \item 我能否用一两句话说明真实目标和最终使用者？
  \item 我是否说明了数据、材料或代码的来源与含义？
  \item 哪些约束对我很明显，但 AI 不可能自动知道？
  \item AI 在这个任务中应当是执行者、副驾驶，还是参谋？
  \item 错误可能造成什么后果，哪些错误最不可接受？
  \item 输出将通过什么独立方法验证？
  \item 是否涉及隐私、保密、版权或组织政策？
  \item 哪些操作必须在执行前由人确认？
  \item 哪些过程和决策需要保留记录？
  \item 最终由谁对结果负责？
\end{enumerate}

如果这些问题大部分无法回答，通常说明任务还没有准备好直接交给 AI。此时，最有价值的第一步不是要求 AI 给出答案，而是让它帮助整理未知信息和澄清问题。

% ----------------------------------------------------------------
\section{实战练习}
% ----------------------------------------------------------------

\subsection*{练习一：判断是否应该使用AI}

从自己的学习、研究或工作中选择三个任务，分别判断：

\begin{enumerate}
  \item 该任务是否适合使用 AI；
  \item AI 应当扮演执行者、副驾驶还是参谋；
  \item 最严重的潜在错误是什么；
  \item 应使用什么方法验证结果。
\end{enumerate}

\subsection*{练习二：为任务选择工具}

针对以下任务，选择合适的 AI 工具类别，并说明理由与风险：

\begin{enumerate}
  \item 理解一个陌生的机器学习概念；
  \item 修改一个包含多个文件的 Python 项目；
  \item 探索一份不含敏感信息的公开交通数据；
  \item 查找支持某项研究结论的原始文献；
  \item 分析包含个人出行轨迹的内部数据；
  \item 为事故应急响应提供决策建议。
\end{enumerate}

\subsection*{练习三：改造一个不可靠的AI任务}

下面的要求看似明确，实际包含许多风险：

\begin{quote}
请自动清洗这份交通数据，选择最好的模型，并告诉我应该采取什么管理措施。
\end{quote}

请将它改造成一个可靠的人机协作流程。至少说明任务需要补充哪些上下文、应拆成哪些步骤、每一步如何验证，以及哪些最终判断不能交给 AI。

\subsection*{练习四：建立个人AI使用规范}

结合自己的任务类型，写一页简短规范，至少包括：

\begin{itemize}
  \item 可以交给 AI 的任务；
  \item 不允许上传的数据；
  \item 必须人工确认的操作；
  \item 必须查阅原始来源的内容；
  \item 需要保存的协作记录；
  \item 发现 AI 输出错误时的处理方式。
\end{itemize}





